\chapter{Родитель равномерного усиления
Callias--\texorpdfstring{$M_4$}{M4} по \texorpdfstring{$H_{15}$}{H15}}
\label{ch:v10-callias-m4-uniform-h15-parent}

Физический пакет распадается на пять типовых блоков
\begin{equation}
 H_{15}=Q_L^{6}\oplus L_L^{2}\oplus u_R^{3}\oplus d_R^{3}\oplus e_R.
\end{equation}
Его три унаследованных заряженных ребра --- $Q_L-u_R$, $Q_L-d_R$ и
$L_L-e_R$. Их incidence-матрица имеет ранг $3$, а лапласиан --- ранг
$3$ и ядро размерности $2$. Ядро состоит из независимых констант на
кварковом компоненте $(Q_L,u_R,d_R)$ и лептонном $(L_L,e_R)$.
После подъёма к пятнадцати каналам их кратности равны $12$ и $3$.

Следовательно, текущий родитель выбирает не один uniform amplifier, а две
независимые амплитуды. Одно дополнительное ребро $Q_L-L_L$ увеличивает
ранг incidence до $4$; дополненный лапласиан имеет ядро
\begin{equation}
 \ker L_{\rm aug}=\operatorname{span}(1,1,1,1,1)^T.
\end{equation}
Его подъём даёт $(1,\ldots,1)^T\in\mathbb C^{15}$ и требуемую Gram-матрицу
$15I_2$ для twist-усилителя.

Но кварк--лептонное ребро отсутствует в унаследованном родителе: его карта
имеет ранг $0$. Даже условное добавление ребра фиксирует только
относительную равномерность; общая амплитуда остаётся свободной.

\begin{theorem}[Двухкомпонентная причина неравномерности]
Три физические связи $H_{15}$ разбивают типовой граф на кварковую и
лептонную компоненты и оставляют две независимые амплитуды. Ровно одно
межкомпонентное ребро условно делает граф связным и выбирает uniform ray.
\end{theorem}

\begin{nogocriterion}[Связность нельзя объявить симметрией]
Полная перестановочная симметрия пятнадцати каналов не унаследована:
заряды и хиральности различают пять блоков. Поэтому uniform vector нельзя
выводить из $S_{15}$-инвариантности; требуется физический кварк--лептонный
коннектор и отдельная нормировка общей амплитуды.
\end{nogocriterion}

\begin{protocol}\begin{enumerate}
\item типовых блоков: \textbf{$5$};
\item унаследованных рёбер: \textbf{$3$};
\item компонент текущего графа: \textbf{$2$};
\item условных дополнительных рёбер: \textbf{$1$};
\item компонент после добавления: \textbf{$1$};
\item inherited bridge rank: \textbf{$0$};
\item physical origin: \textbf{$0/3$};
\item следующий гейт: \textbf{аудит кварк--лептонного коннектора}.
\end{enumerate}\end{protocol}

Машинный сертификат сохранён в
\path{s2t_v10_particle_wrinkle_dislocation_callias_equal_charge_twist_m4_uniform_h15_amplification_parent_origin_gate_results.json}.